\documentclass[11pt,a4paper]{article}
\usepackage[T1]{fontenc}
\usepackage[utf8]{inputenc}
\usepackage[french]{babel}
\usepackage{lmodern}
\usepackage{textcomp}
\usepackage{amsmath,amssymb}
\usepackage{geometry}
\usepackage{tabularx,array}
\usepackage[table]{xcolor}
\usepackage{tikz}
\geometry{left=8mm,right=8mm,top=7mm,bottom=7mm}
\setlength{\parindent}{0pt}
\pagestyle{empty}
\setlength{\tabcolsep}{3pt}
\renewcommand{\arraystretch}{1.05}
\renewcommand{\tabularxcolumn}[1]{m{#1}}

% --- Palette pensée pour une impression noir et blanc ---
\definecolor{canlight}{gray}{0.93}
\definecolor{canmid}{gray}{0.72}
\definecolor{candark}{gray}{0.22}
\definecolor{answerblue}{RGB}{0,82,165}

% --- Colonnes ---
\newcolumntype{Q}{>{\centering\arraybackslash}X}
\newcolumntype{N}{>{\centering\arraybackslash}m{6mm}}
\newcolumntype{R}{>{\centering\arraybackslash}m{22mm}}

% --- En-tête ---
\newcommand{\canheader}[1]{%
  \begin{center}
    {\LARGE\bfseries COURSE AUX NOMBRES}\par
    \vspace{0.7mm}
    {\color{canmid}\rule{27mm}{0.45pt}\hspace{2mm}\color{black}$\diamond$\hspace{1.4mm}{\large\bfseries #1}\hspace{1.4mm}$\diamond$\hspace{2mm}\color{canmid}\rule{27mm}{0.45pt}\color{black}}
  \end{center}
  \vspace{0.8mm}
  \renewcommand{\arraystretch}{1.25}%
  \begin{tabularx}{\linewidth}{>{\centering\arraybackslash}X >{\centering\arraybackslash}X >{\centering\arraybackslash}m{40mm}}
    \rowcolor{canlight}
    \footnotesize\textbf{30 questions} -- 9 minutes &
    \footnotesize\textbf{Sans calculatrice} &
    \footnotesize\textbf{Score :} \rule{13mm}{0.45pt}\,/30
  \end{tabularx}
  \renewcommand{\arraystretch}{1.05}%
  \vspace{1mm}
  \begin{center}
    {\scriptsize Réponse uniquement \quad$\diamond$\quad Reporter le score sur la fiche annuelle.}
  \end{center}
  \vspace{1.2mm}
}

\newcommand{\qrow}[2]{\textbf{#1} & \parbox[c][15.0mm][c]{\linewidth}{\centering #2} & \\ \hline}
\newcommand{\qrowfig}[2]{\textbf{#1} & \parbox[c][27mm][c]{\linewidth}{\centering #2} & \\ \hline}
\newcommand{\qrowcorr}[3]{\textbf{#1} & \parbox[c][15.0mm][c]{\linewidth}{\centering #2} & {\color{answerblue}\bfseries #3} \\ \hline}
\newcommand{\qrowfigcorr}[3]{\textbf{#1} & \parbox[c][27mm][c]{\linewidth}{\centering #2} & {\color{answerblue}\bfseries #3} \\ \hline}

\newcommand{\cantable}[1]{%
\arrayrulecolor{canmid}
\begin{tabularx}{\linewidth}{N|Q|R}
\rowcolor{canlight}\textbf{N\textsuperscript{o}}&\centering\textbf{Question}&\textbf{Réponse}\\\hline
#1
\end{tabularx}
\arrayrulecolor{black}
}

\newcommand{\hypotfigure}{%
\begin{tikzpicture}[baseline=(current bounding box.center),scale=0.72,line width=0.6pt]
  \node[font=\bfseries\normalsize] at (1.6,2.55) {Quelle est l'hypoténuse ?};
  \coordinate (A) at (0,0); \coordinate (B) at (0,1.85); \coordinate (C) at (3.2,0);
  \draw (A)--(B)--(C)--cycle;
  \draw (0,0) rectangle (0.32,0.32);
  \node[below left] at (A) {\small A};
  \node[above left] at (B) {\small B};
  \node[below right] at (C) {\small C};
\end{tikzpicture}}


\newcommand{\canheadercorr}[1]{%
  \begin{center}
    {\LARGE\bfseries COURSE AUX NOMBRES}\par
    \vspace{0.7mm}
    {\color{canmid}\rule{27mm}{0.45pt}\hspace{2mm}\color{black}$\diamond$\hspace{1.4mm}{\large\bfseries #1}\hspace{1.4mm}$\diamond$\hspace{2mm}\color{canmid}\rule{27mm}{0.45pt}\color{black}}
  \end{center}
  \vspace{0.8mm}
  \renewcommand{\arraystretch}{1.25}%
  \begin{tabularx}{\linewidth}{>{\centering\arraybackslash}X >{\centering\arraybackslash}X >{\centering\arraybackslash}m{40mm}}
    \rowcolor{canlight}
    \footnotesize\textbf{30 questions} -- 9 minutes &
    \footnotesize\textbf{Sans calculatrice} &
    \footnotesize\textcolor{answerblue}{\textbf{CORRIGÉ}}
  \end{tabularx}
  \renewcommand{\arraystretch}{1.05}%
  \vspace{1mm}
  \begin{center}
    {\scriptsize Les réponses apparaissent en \textcolor{answerblue}{\textbf{bleu}}.}
  \end{center}
  \vspace{1.2mm}
}


\begin{document}
\canheader{CAN 5}

\begin{minipage}[t]{0.493\textwidth}
\normalsize
\cantable{
\qrow{1}{$12\times7$}
\qrow{2}{$800\div20$}
\qrow{3}{$-7+5-3+1$}
\qrow{4}{Compléter : $105-\ldots=89$}
\qrow{5}{$20\,\%$ de 45}
\qrow{6}{Le triple de la somme de 4 et 5}
\qrow{7}{$888-199$}
\qrow{8}{Quel est le nombre qui, multiplié par 4, donne 7 ? Donner une fraction.}
\qrow{9}{$0{,}0009\times100$}
\qrow{10}{Médiane de la série : 3 ; 8 ; 5 ; 7 ; 6}
\qrow{11}{Le plus grand : $\dfrac23$ ou $\dfrac35$ ?}
\qrow{12}{L'opposé de $-11$}
\qrow{13}{$\dfrac58$ de 32}
\qrow{14}{4 h 12 min = $\ldots$ min}
\qrow{15}{$8-8\div2+6$}
}
\end{minipage}\hfill
\begin{minipage}[t]{0.493\textwidth}
\normalsize
\cantable{
\qrow{16}{$\dfrac35+\dfrac15$}
\qrow{17}{Aire d'un carré de côté 7 cm}
\qrow{18}{Dans une série de 40 valeurs, une valeur apparaît 8 fois. Sa fréquence décimale est :}
\qrow{19}{Écrire $1{,}75$ sous la forme d'une fraction de dénominateur 4.}
\qrow{20}{Un cube de 1 cm de côté a un volume de :}
\qrow{21}{Quel est le plus petit multiple de 7 strictement supérieur à 30 ?}
\qrow{22}{$1{,}8\times98+2\times1{,}8$}
\qrow{23}{$1\,\mathrm{m}^3=\ldots$ L}
\qrow{24}{4 compas coûtent 10 \texteuro. Combien coûtent 6 compas ?}
\qrow{25}{Dans un triangle rectangle, un côté mesure 8 cm. L'hypoténuse peut-elle mesurer 5 cm ?}
\qrow{26}{Quel nombre, multiplié par lui-même, donne 100 ?}
\qrow{27}{$49{,}8+51{,}1$ est plus proche de 10, de 100 ou de 1 000 ?}
\qrow{28}{Encadrer $12{,}34$ par deux entiers consécutifs.}
\qrow{29}{Réduire : $5x-2x$}
\qrow{30}{$3x+2$ est-il une somme ou un produit ?}
}
\end{minipage}

\clearpage

\canheadercorr{CAN 5}

\begin{minipage}[t]{0.493\textwidth}
\normalsize
\cantable{
\qrowcorr{1}{$12\times7$}{$84$}
\qrowcorr{2}{$800\div20$}{$40$}
\qrowcorr{3}{$-7+5-3+1$}{$-4$}
\qrowcorr{4}{Compléter : $105-\ldots=89$}{$16$}
\qrowcorr{5}{$20\,\%$ de 45}{$9$}
\qrowcorr{6}{Le triple de la somme de 4 et 5}{$27$}
\qrowcorr{7}{$888-199$}{$689$}
\qrowcorr{8}{Quel est le nombre qui, multiplié par 4, donne 7 ? Donner une fraction.}{$\dfrac74$}
\qrowcorr{9}{$0{,}0009\times100$}{$0{,}09$}
\qrowcorr{10}{Médiane de la série : 3 ; 8 ; 5 ; 7 ; 6}{$6$}
\qrowcorr{11}{Le plus grand : $\dfrac23$ ou $\dfrac35$ ?}{$\dfrac23$}
\qrowcorr{12}{L'opposé de $-11$}{$11$}
\qrowcorr{13}{$\dfrac58$ de 32}{$20$}
\qrowcorr{14}{4 h 12 min = $\ldots$ min}{252 min}
\qrowcorr{15}{$8-8\div2+6$}{$10$}
}
\end{minipage}\hfill
\begin{minipage}[t]{0.493\textwidth}
\normalsize
\cantable{
\qrowcorr{16}{$\dfrac35+\dfrac15$}{$\dfrac45$}
\qrowcorr{17}{Aire d'un carré de côté 7 cm}{49 cm$^2$}
\qrowcorr{18}{Dans une série de 40 valeurs, une valeur apparaît 8 fois. Sa fréquence décimale est :}{$0{,}2$}
\qrowcorr{19}{Écrire $1{,}75$ sous la forme d'une fraction de dénominateur 4.}{$\dfrac74$}
\qrowcorr{20}{Un cube de 1 cm de côté a un volume de :}{1 cm$^3$}
\qrowcorr{21}{Quel est le plus petit multiple de 7 strictement supérieur à 30 ?}{$35$}
\qrowcorr{22}{$1{,}8\times98+2\times1{,}8$}{$180$}
\qrowcorr{23}{$1\,\mathrm{m}^3=\ldots$ L}{$1000$ L}
\qrowcorr{24}{4 compas coûtent 10 \texteuro. Combien coûtent 6 compas ?}{15 \texteuro}
\qrowcorr{25}{Dans un triangle rectangle, un côté mesure 8 cm. L'hypoténuse peut-elle mesurer 5 cm ?}{non}
\qrowcorr{26}{Quel nombre, multiplié par lui-même, donne 100 ?}{$10$}
\qrowcorr{27}{$49{,}8+51{,}1$ est plus proche de 10, de 100 ou de 1 000 ?}{$100$}
\qrowcorr{28}{Encadrer $12{,}34$ par deux entiers consécutifs.}{$12<12{,}34<13$}
\qrowcorr{29}{Réduire : $5x-2x$}{$3x$}
\qrowcorr{30}{$3x+2$ est-il une somme ou un produit ?}{somme}
}
\end{minipage}
\end{document}
